% GENERATED FILE. Edit canonical lessons, then run npm run generate:books.
% canonical-source-hash: fnv1a64:51c98ad451bf83b0
% canonical-lessons: KA-C362-nagarikate, KA-C362-sthapisu, KA-C362-nirmisu, KA-C362-kedavu, KA-C362-pratime

\chapter{Builders of Old}
\label{ch:ka-builders-of-old}
\hlchaptermodality{362}

\begin{chapteropening}
\textbf{By the end of this chapter:} \emph{I can say a civilization, to found, to construct, to demolish and a statue.}

Complete the last lesson of chapter 362: I can say a civilization, to found, to construct, to demolish and a statue.
\end{chapteropening}

\section[\texorpdfstring{\kn{ನಾಗರಿಕತೆ} (nāgarikate)}{ನಾಗರಿಕತೆ (nāgarikate)}]{\kn{ನಾಗರಿಕತೆ} (nāgarikate) — a civilization}

\label{lesson:KA-C362-nagarikate}



\begin{quote}
Before the new one: say the Kannada for a bastion, then the Kannada for to bury.
\end{quote}

\subsection*{You'll want to know: \kn{ನಾಗರಿಕತೆ}}
\textbf{\kn{ನಾಗರಿಕತೆ}} — \emph{nāgarikate} — \textquotedblleft{}a civilization\textquotedblright{}.

\begin{quote}
\kn{ನದಿಯ} \kn{ದಡದಲ್ಲಿ} \kn{ಒಂದು} \kn{ಪ್ರಾಚೀನ} \kn{ನಾಗರಿಕತೆ} \kn{ಬೆಳೆಯಿತು}.
\end{quote}

\begin{quote}
\emph{nadiya daḍadalli ondu prācīna nāgarikate beḷeyitu.}
\end{quote}

\begin{quote}
An ancient civilization grew up on the river bank.
\end{quote}

Plural: \kn{ನಾಗರಿಕತೆಗಳು}, \emph{nāgarikategaḷu}.

-\kn{ತೆ} makes an abstract noun of an adjective: town-bred ways become civilization.

\begin{cousinweb}
\textbf{\kn{ನಾಗರಿಕತೆ}} (\emph{nāgarikate}) is built on Sanskrit \emph{nāgarika}, \textquotedblleft{}born or living in a town, civic, polite\textquotedblright{} (Monier-Williams), with the Kannada abstract ending -\emph{te}.
\end{cousinweb}

\begin{grammarlens}[title={What you've built}]
One more word for telling what happened.
\end{grammarlens}

\subsection*{Your turn}
\begin{itemize}
\raggedright
  \item \emph{Say it:} \emph{nāgarikate}
  \item \emph{Say it:} \emph{nāgarikate}, once more
  \item \emph{Say it:} \emph{nadiya daḍadalli ondu prācīna nāgarikate beḷeyitu.} — An ancient civilization grew up on the river bank.
  \item \emph{Recall:} say the Kannada for a bastion, then the Kannada for to bury, then say \emph{nāgarikate} again
\end{itemize}

\begin{tcolorbox}[breakable,colback=teal!4,colframe=teal!35!black,title={Before you move on}]
What does \kn{ನಾಗರಿಕತೆ} mean? (\textquotedblleft{}A civilization\textquotedblright{}.) Say it once more.
\end{tcolorbox}
\section[\texorpdfstring{\kn{ಸ್ಥಾಪಿಸು} (sthāpisu)}{ಸ್ಥಾಪಿಸು (sthāpisu)}]{\kn{ಸ್ಥಾಪಿಸು} (sthāpisu) — to found; to establish}

\label{lesson:KA-C362-sthapisu}



\begin{quote}
Before the new one: say the Kannada for to bury, then the Kannada for a civilization.
\end{quote}

\subsection*{You'll want to know: \kn{ಸ್ಥಾಪಿಸು}}
\textbf{\kn{ಸ್ಥಾಪಿಸು}} — \emph{sthāpisu} — \textquotedblleft{}to found; to establish\textquotedblright{}.

\begin{quote}
\kn{ಆ} \kn{ರಾಜ} \kn{ಈ} \kn{ನಗರವನ್ನು} \kn{ಐನೂರು} \kn{ವರ್ಷದ} \kn{ಹಿಂದೆ} \kn{ಸ್ಥಾಪಿಸಿದನು}.
\end{quote}

\begin{quote}
\emph{ā rāja ī nagaravannu ainūru varṣada hinde sthāpisidanu.}
\end{quote}

\begin{quote}
That king founded this city five hundred years ago.
\end{quote}

Past (I …): \kn{ಸ್ಥಾಪಿಸಿದೆ}, \emph{sthāpiside}.

Cities, schools and companies are founded: \kn{ಶಾಲೆಯನ್ನು} \kn{ಸ್ಥಾಪಿಸು}.

\begin{cousinweb}
\textbf{\kn{ಸ್ಥಾಪಿಸು}} (\emph{sthāpisu}) is built on Sanskrit \emph{sthāpana}, \textquotedblleft{}causing to stand, fixing\textquotedblright{} (CDIAL 13755), from \emph{sthāpayati}, \textquotedblleft{}makes stand\textquotedblright{}, with the Kannada verb ending -\emph{isu}.
\end{cousinweb}

\begin{grammarlens}[title={What you've built}]
One more word for telling what happened.
\end{grammarlens}

\subsection*{Your turn}
\begin{itemize}
\raggedright
  \item \emph{Say it:} \emph{sthāpisu}
  \item \emph{Say it:} \emph{sthāpisu}, once more
  \item \emph{Say it:} \emph{ā rāja ī nagaravannu ainūru varṣada hinde sthāpisidanu.} — That king founded this city five hundred years ago.
  \item \emph{Recall:} say the Kannada for to bury, then the Kannada for a civilization, then say \emph{sthāpisu} again
\end{itemize}

\begin{tcolorbox}[breakable,colback=teal!4,colframe=teal!35!black,title={Before you move on}]
What does \kn{ಸ್ಥಾಪಿಸು} mean? (\textquotedblleft{}To found; to establish\textquotedblright{}.) Say it once more.
\end{tcolorbox}
\section[\texorpdfstring{\kn{ನಿರ್ಮಿಸು} (nirmisu)}{ನಿರ್ಮಿಸು (nirmisu)}]{\kn{ನಿರ್ಮಿಸು} (nirmisu) — to construct; to build}

\label{lesson:KA-C362-nirmisu}



\begin{quote}
Before the new one: say the Kannada for a civilization, then the Kannada for to found.
\end{quote}

\subsection*{You'll want to know: \kn{ನಿರ್ಮಿಸು}}
\textbf{\kn{ನಿರ್ಮಿಸು}} — \emph{nirmisu} — \textquotedblleft{}to construct; to build\textquotedblright{}.

\begin{quote}
\kn{ಕೆಲಸಗಾರರು} \kn{ನದಿಗೆ} \kn{ಹೊಸ} \kn{ಸೇತುವೆ} \kn{ನಿರ್ಮಿಸಿದರು}.
\end{quote}

\begin{quote}
\emph{kelasagāraru nadige hosa sētuve nirmisidaru.}
\end{quote}

\begin{quote}
The workers built a new bridge over the river.
\end{quote}

Past (I …): \kn{ನಿರ್ಮಿಸಿದೆ}, \emph{nirmiside}.

A formal partner of \kn{ಕಟ್ಟು}. Films are \textquotedblleft{}made\textquotedblright{} with it too: \kn{ಸಿನಿಮಾ} \kn{ನಿರ್ಮಿಸು}.

\begin{cousinweb}
\textbf{\kn{ನಿರ್ಮಿಸು}} (\emph{nirmisu}) is built on Sanskrit \emph{nirmāṇa}, \textquotedblleft{}measuring, making, creating\textquotedblright{} (CDIAL 7368), with the Kannada verb ending -\emph{isu}.
\end{cousinweb}

\begin{grammarlens}[title={What you've built}]
One more word for telling what happened.
\end{grammarlens}

\subsection*{Your turn}
\begin{itemize}
\raggedright
  \item \emph{Say it:} \emph{nirmisu}
  \item \emph{Say it:} \emph{nirmisu}, once more
  \item \emph{Say it:} \emph{kelasagāraru nadige hosa sētuve nirmisidaru.} — The workers built a new bridge over the river.
  \item \emph{Recall:} say the Kannada for a civilization, then the Kannada for to found, then say \emph{nirmisu} again
\end{itemize}

\begin{tcolorbox}[breakable,colback=teal!4,colframe=teal!35!black,title={Before you move on}]
What does \kn{ನಿರ್ಮಿಸು} mean? (\textquotedblleft{}To construct; to build\textquotedblright{}.) Say it once more.
\end{tcolorbox}
\section[\texorpdfstring{\kn{ಕೆಡವು} (keḍavu)}{ಕೆಡವು (keḍavu)}]{\kn{ಕೆಡವು} (keḍavu) — to demolish; to knock down}

\label{lesson:KA-C362-kedavu}



\begin{quote}
Before the new one: say the Kannada for to found, then the Kannada for to construct.
\end{quote}

\subsection*{You'll want to know: \kn{ಕೆಡವು}}
\textbf{\kn{ಕೆಡವು}} — \emph{keḍavu} — \textquotedblleft{}to demolish; to knock down\textquotedblright{}.

\begin{quote}
\kn{ಅವರು} \kn{ಹಳೆಯ} \kn{ಕಟ್ಟಡವನ್ನು} \kn{ಕೆಡವಿದರು}.
\end{quote}

\begin{quote}
\emph{avaru haḷeya kaṭṭaḍavannu keḍavidaru.}
\end{quote}

\begin{quote}
They demolished the old building.
\end{quote}

Past (I …): \kn{ಕೆಡವಿದೆ}, \emph{keḍavide}.

What is knocked down takes -\kn{ಅನ್ನು}. A wrestler can \kn{ಕೆಡವು} his rival too.

\begin{cousinweb}
\textbf{\kn{ಕೆಡವು}} (\emph{keḍavu}) is native Dravidian: DEDR 1524 lists Kannada \emph{keḍavu}, \textquotedblleft{}to throw down, to fell\textquotedblright{}, under words for lying down.
\end{cousinweb}

\begin{grammarlens}[title={What you've built}]
One more word for telling what happened.
\end{grammarlens}

\subsection*{Your turn}
\begin{itemize}
\raggedright
  \item \emph{Say it:} \emph{keḍavu}
  \item \emph{Say it:} \emph{keḍavu}, once more
  \item \emph{Say it:} \emph{avaru haḷeya kaṭṭaḍavannu keḍavidaru.} — They demolished the old building.
  \item \emph{Recall:} say the Kannada for to found, then the Kannada for to construct, then say \emph{keḍavu} again
\end{itemize}

\begin{tcolorbox}[breakable,colback=teal!4,colframe=teal!35!black,title={Before you move on}]
What does \kn{ಕೆಡವು} mean? (\textquotedblleft{}To demolish; to knock down\textquotedblright{}.) Say it once more.
\end{tcolorbox}
\section[\texorpdfstring{\kn{ಪ್ರತಿಮೆ} (pratime)}{ಪ್ರತಿಮೆ (pratime)}]{\kn{ಪ್ರತಿಮೆ} (pratime) — a statue}

\label{lesson:KA-C362-pratime}



\begin{quote}
Before the new one: say the Kannada for to construct, then the Kannada for to demolish.
\end{quote}

\subsection*{You'll want to know: \kn{ಪ್ರತಿಮೆ}}
\textbf{\kn{ಪ್ರತಿಮೆ}} — \emph{pratime} — \textquotedblleft{}a statue\textquotedblright{}.

\begin{quote}
\kn{ವೃತ್ತದಲ್ಲಿ} \kn{ಆ} \kn{ರಾಜನ} \kn{ದೊಡ್ಡ} \kn{ಪ್ರತಿಮೆ} \kn{ಇದೆ}.
\end{quote}

\begin{quote}
\emph{vṛttadalli ā rājana doḍḍa pratime ide.}
\end{quote}

\begin{quote}
There is a large statue of that king at the circle.
\end{quote}

Plural: \kn{ಪ್ರತಿಮೆಗಳು}, \emph{pratimegaḷu}.

\begin{cousinweb}
\textbf{\kn{ಪ್ರತಿಮೆ}} (\emph{pratime}) is borrowed from Sanskrit \emph{pratimā}, \textquotedblleft{}an image, a likeness, a statue\textquotedblright{} (Monier-Williams).
\end{cousinweb}

\begin{grammarlens}[title={What you've built}]
One more word for telling what happened.
\end{grammarlens}

\subsection*{Your turn}
\begin{itemize}
\raggedright
  \item \emph{Say it:} \emph{pratime}
  \item \emph{Say it:} \emph{pratime}, once more
  \item \emph{Say it:} \emph{vṛttadalli ā rājana doḍḍa pratime ide.} — There is a large statue of that king at the circle.
  \item \emph{Recall:} say the Kannada for to construct, then the Kannada for to demolish, then say \emph{pratime} again
\end{itemize}

\begin{tcolorbox}[breakable,colback=teal!4,colframe=teal!35!black,title={Before you move on}]
What does \kn{ಪ್ರತಿಮೆ} mean? (\textquotedblleft{}A statue\textquotedblright{}.) Say it once more.
\end{tcolorbox}
